\section{The Adaptive Tree of Experts (AToE) PINN method}
\label{sec:method}

\subsection{Overview of the Method}
\label{sec:overview}

The tree based structure described in
\S\ref{sec:background} forms a natural starting point for a wide range of domain decomposition
designs that combine local models into a Mixture of Experts (MoE). In this
section we present such a method which we refer to as
AToE (Adaptive Tree of Experts). Section~\ref{sec:additional_methods}
describes additional variants of the main AToE algorithm that can further improve its performance. The full procedure is summarized in Algorithm~\ref{alg:atoe}, where each of
its stages is described in detail in the subsections that follow.

\begin{algorithm}[H]
\caption{AToE: Adaptive Tree of Experts}
\label{alg:atoe}
\begin{algorithmic}[1]
\Require PDE problem, template NN architecture $\mathcal{A}_{\mathrm{base}}$, template expert sub-net architecture $\mathcal{A}_{\mathrm{exp}}$, budget of sub-network local experts $M$
\State Train a vanilla PINN $u_0$ based on the $\mathcal{A}_{\mathrm{base}}$ architecture
\State Build an approximating tree to $u_0$; trim to an $M$-term complete binary tree
\State Record $u_0$ on the interfaces of the tree's leaves $\{\Omega_j\}$, $\Omega\times[0,T]=\cup_j \Omega_j$
\ForAll{leaf regions $\Omega_j$ \textbf{in parallel}}
    \State Train an expert $u_j$ with architecture  $\mathcal{A}_{\mathrm{exp}}$ on $\Omega_j$, using the values of $u_0$ as sub-region boundary conditions
\EndFor
\State Blend the experts via smoothstep POU (partition of unity)
\State Perform fine-tune training of the blended composition under the global PINN loss
\Ensure Global approximate solution $u_\theta$
\end{algorithmic}
\end{algorithm}

\vspace{0.5em}

\subsection{Learning the Initial Network}
\label{sec:root_training}

The goal of this initial step is to obtain a coarse approximation of the PDE solution that is sufficient to detect local spatiotemporal structures of the true solution.

The architecture and training process of the initial model $u_0$ follows the classic PINN formulation described in
\S\ref{sec:background}. For the second order benchmarks we use an MLP with 4 hidden layers of 60 neurons. Training consists of 1k Adam iterations
with learning rate $10^{-3}$, followed by 40k SSBroyden iterations, and all computations are performed in
double precision, as suggested in~\cite{kiyani2025optimizer}.

For the higher order benchmarks we use larger networks: an MLP with 6 hidden layers of 128 neurons, combined with a random Fourier feature embedding~\cite{tancik2020fourier} of the input, with scale $2$ and embedding dimension equal to the layer width. To keep a network of this size affordable we train it in single precision, and the reduced precision in turn calls for a longer first order stage: 30k Adam iterations governed by a scheduler with base learning rate $10^{-3}$, a warmup from $10^{-5}$ to $10^{-3}$ over the first 2k iterations, and an exponential decay by a factor of $0.9$ every 2k iterations thereafter, followed by 30k L-BFGS iterations with learning rate $0.1$.

For KS, whose chaotic dynamics make a single global fit over the full horizon difficult, the root itself is trained by the sequential time marching approach of~\cite{wight2021timeadaptive}. Advancing over successive time windows, rather than fitting the full horizon at once, is the common treatment of KS in the PINN literature. The time interval is divided into 5 equal windows, and an independent network with the higher order architecture above is trained on each window in order, with the converged prediction of the preceding network at the shared time slice serving as the initial condition of the next. The resulting piecewise model over the 5 windows serves as the root $u_0$, and the remainder of the pipeline, the tree construction with time tiling, the expert training, and the global fine-tuning, proceeds exactly as for KdV.

In both regimes, all stages use full batch training with 20k residual points and 256 points each for the initial and boundary conditions; the residual set is redrawn at every Adam iteration and held fixed during the quasi-Newton stages. The loss weights are $\lambda_r = \lambda_{ic} = \lambda_{bc} = 1$ for the second order benchmarks, while the higher order ones use $\lambda_{ic} = \lambda_{bc} = 1000$ with $\lambda_r = 1$; these weights are kept unchanged across all training phases.

Typically, the initial approximation $u_0$ has large errors in spatiotemporal regions where the solution develops transients and shocks (see for example Figure \ref{fig:results_burgers}(c) where a vanilla PINN has large errors around $x=0$ at times $0.3\le t \le 1$ for the Burgers equation). 

\vspace{0.5em}

\subsection{Construction of the Approximating Tree and Domain Decomposition}
\label{sec:tree_construction}

The goal of the decomposition is to partition the domain into regions that
reflect the local complexity of the solution. A geometric wavelet tree
$\mathcal{T}$ is constructed from the initial approximation $u_0$ and
compressed using an $M$-term approximation \eqref{m-term}. To preserve a clean
hierarchical structure, both the ancestors and the siblings of the selected nodes are also retained, so that the accepted regions form a complete binary sub-tree (that is not necessarily balanced).

For the higher order benchmarks the tree is not built on the full space--time domain directly. The solutions of these equations develop
localized features that move and evolve in time, so a small number of
axis-aligned splits cannot efficiently isolate them over the full time horizon: a spatial subdomain that is simple at one time contains a feature at another. To keep each expert region internally simple, we tile the time axis into equal-width tiles and build an independent tree inside each tile. A second motivation is that the tiling imposes a structure similar to sequential time-marching~\cite{wight2021timeadaptive}: each expert's local initial condition originates at most one tile width in the past, keeping every local problem anchored to temporally nearby data.

The global budget $M$ is distributed across the tiles according to a prescribed allocation strategy
(e.g., uniform, linear, or quadratic), and an independent $M_i$-term
approximation is constructed within each tile, with $\sum_i M_i = M$. The
leaves of all tile trees together form the decomposition, and the rest of the
pipeline proceeds unchanged.

Each leaf region $\Omega_j \in \Lambda$ of the resulting tree is assigned a neural network expert $u_j$ with the expert architecture
$\mathcal{A}_{\mathrm{exp}}$, which is smaller than the root architecture. The reduced size is possible without sacrificing accuracy because each expert approximates the solution only on its own subdomain, where by construction of the tree the solution has more uniform local regularity, rather than the full problem with all its variation. 

Each expert is an MLP with 3 hidden layers of 24 neurons for the second order benchmarks, and 3 hidden layers of 30 neurons, with the random Fourier feature embedding of Section~\ref{sec:root_training}, for the higher order ones. In the experiments shown in Section~\ref{sec:experiments} we set $M = 8$ for the second order PDEs and $M = 18$ for the higher order ones, using 3 time tiles, with a linear allocation strategy for KdV and a quadratic one for KS.
The KS solution develops fine structure that intensifies toward the end of
the time horizon more sharply than the KdV solution does, so its budget is
skewed more strongly toward the later tiles. Since the completion to a full binary tree adds or removes leaves relative to the selected terms, the resulting number of experts varies per problem (see Tables~\ref{tab:benchmark_results_second} and~\ref{tab:benchmark_results_high}).

\subsection{Interpretation of the Adaptive Decomposition}
\label{sec:interpretation}

As discussed in \S\ref{sec:background}, the equivalence between the
Tree-Besov semi-norm and the geometric wavelet norm implies that the decay of
geometric wavelet coefficients is directly related to the smoothness of the
underlying function.

A large geometric wavelet norm at a single node does not necessarily
imply that the corresponding region is non-smooth, since nodes near the root
of the tree may simply capture large scale variations of the solution.
Instead, smoothness is reflected in how rapidly the geometric wavelet norms
decay along a branch of the tree. Regions associated with branches whose
coefficients remain significant at increasing depths exhibit slower decay and
therefore correspond to lower regularity and higher local complexity.

Viewed through these lens, the tree construction procedure described
in \S\ref{sec:tree_construction} naturally favors the retention of
branches that exhibit slow coefficient decay across multiple levels of the
hierarchy. Consequently, the resulting decomposition concentrates refinement
in the less smooth and more complex portions of the domain, while maintaining
a coarser representation in regions where the coefficient decay is more rapid.

Extended methods that take further advantage of our theoretical setup are detailed in \S\ref{sec:additional_methods}. However, even the
simple choice of a fixed expert architecture already applies a form of
adaptive capacity allocation. Regions associated with branches that persist to
greater depths receive a larger number of experts and therefore a larger
effective approximation capacity, whereas smoother regions are represented by
shallower portions of the tree. As illustrated in
Figure~\ref{fig:capacity_maps}, this mechanism produces a nonlinear
distribution of capacity across the domain that reflects the underlying
complexity of the solution.

\vspace{0.5em}

\begin{figure}[H]
\centering

\begin{subfigure}{0.45\textwidth}
\centering
\includegraphics[width=\textwidth]
{images/section3_method/schrodinger_capacity_map.png}
\caption{Schrödinger}
\label{fig:schrodinger_capacity}
\end{subfigure}
\hfill
\begin{subfigure}{0.45\textwidth}
\centering
\includegraphics[width=\textwidth]
{images/section3_method/allen_cahn_capacity_map.png}
\caption{Allen--Cahn}
\label{fig:allen_cahn_capacity}
\end{subfigure}

\vspace{0.5em}

\begin{subfigure}{0.45\textwidth}
\centering
\includegraphics[width=\textwidth]
{images/section3_method/ks_capacity_map.png}
\caption{KS}
\end{subfigure}
\hfill
\begin{subfigure}{0.45\textwidth}
\centering
\includegraphics[width=\textwidth]
{images/section3_method/kdv_capacity_map.png}
\caption{KdV}
\end{subfigure}

\caption{
Effective approximation capacity induced by the adaptive decomposition relative to the subdomain volume for benchmark PDEs. Trees are constructed using the $M$-term procedure of Section~\ref{sec:tree_construction} with $M=8$ for the second order examples and $M=18$ with time tiling for the higher order ones.}
\label{fig:capacity_maps}
\end{figure}

\vspace{0.5em}

\subsection{Expert Training}
\label{sec:expert_training}

Before we assemble and fine-tune the experts together through the global loss terms, we exploit a
core MoE idea and let each expert solve a localized, and therefore easier problem on its own subdomain. Inspired by the progressive scheme of~\cite{luo2025pdd}, each local problem is made independent by using evaluations of $u_0$ and its partial derivatives as boundary conditions on the shared interfaces. Concretely, on each leaf subdomain $\Omega_j$ we minimize a local
objective
\begin{equation}
\mathcal{L}_j
= \lambda_r\,\mathcal{L}^{\mathrm{res}}_j
+ \lambda_{ic}\,\mathcal{L}^{\mathrm{ic}}_j
+ \lambda_{bc}\,\mathcal{L}^{\mathrm{bc}}_j
+ \lambda_{\Gamma}\,\mathcal{L}^{\Gamma}_j,
\label{eq:leaf_local_loss}
\end{equation}
where $\mathcal{L}^{\mathrm{res}}_j$ enforces the PDE residual on collocation
points inside $\Omega_j$, and $\mathcal{L}^{\mathrm{ic}}_j$,
$\mathcal{L}^{\mathrm{bc}}_j$ enforce the initial condition $f_0$ and given boundary conditions
on the parts of $\partial\Omega_j$ that coincide with the boundary of
$\Omega$.

The region $\Omega_j$ generally borders a different neighbor along
each spatial dimension, so let $\mathcal{I}_j \subseteq \{1,\ldots,d\}$
denote the dimensions along which $\Omega_j$ has an interior neighbor, as
opposed to coinciding with the physical boundary of $\Omega$, and for
$i \in \mathcal{I}_j$ let $X^{\Gamma_j}_i$ denote the sample points on the
corresponding interior face of $\Omega_j$. Let
$x=(x_1,\ldots,x_d)\in\mathbb{R}^d$ denote the spatial coordinates and $N$
the differential order of the PDE. On the face normal to $x_i$ we match the
value together with the derivatives of orders $m=1,\ldots,N-1$.

Matching the value alone would pin down each expert only up to a mismatch in its
derivatives at the interface, which the PDE operator differentiates further. Matching derivatives up to order $N-1$ encourages approximate $C^{N-1}$
agreement between neighboring experts, the level at which local solutions
of an order-$N$ equation compose into a global one. The agreement is only
approximate, and it is the smooth blending of
\S\ref{sec:smoothstep_windows} that yields a globally smooth model. Since all interface data are generated by the single approximate solution $u_0$, these conditions are mutually consistent by construction.

Writing $X^{\Gamma}_j = \bigcup_{i\in\mathcal{I}_j} X^{\Gamma_j}_i$ for the
union over all such faces, the loss ties $u_j$ and its derivatives, on every
interior face, to the frozen initial prediction $u_0$:
\begin{equation}
\mathcal{L}^{\Gamma}_j
= \frac{1}{N\,|X^{\Gamma}_j|}
\sum_{i\in\mathcal{I}_j}
\sum_{(x,t)\in X^{\Gamma_j}_i}
\sum_{m=0}^{N-1}
\left| \frac{\partial^m u_j}{\partial^m x_i}(x,t)
- \frac{\partial^m u_0}{\partial^m x_i}(x,t) \right|^2 .
\label{eq:interface_loss}
\end{equation}
In this way every expert solves a localized problem whose interface data on
the interior interfaces is supplied by the root, so the leaves can be trained independently, and in parallel, while remaining mutually consistent. 

For the higher order benchmarks, two practical adjustments are made to the derivative matching in \eqref{eq:interface_loss}. In single precision, and with solutions whose derivative magnitudes grow steeply with order, the mismatch terms of orders $m \geq 2$ attain values that dominate the gradient budget and destabilize training. We therefore truncate the matching to the value and the first derivative, and normalize the first derivative term of each expert by the scale
\begin{equation}
s_1 = 1 + \frac{1}{|X^{\Gamma}_j|}\sum_{i\in\mathcal{I}_j}\sum_{(x,t)\in X^{\Gamma_j}_i} \left|\frac{\partial u_0}{\partial x_i}(x,t)\right|^2,
\label{eq:derivative_scale}
\end{equation}
computed over the expert's interior spatial faces, so that the derivative mismatch becomes a relative quantity of magnitude comparable to the value term. The $+1$ floor guarantees that normalization can only tame a large term, never amplify a small one.

The experts are trained with the same optimizer
schedule as the root of the corresponding benchmark (Section~\ref{sec:root_training}). Each leaf draws a fixed, uniformly sampled point set whose size is proportional to its volume, subject to a
minimum count, so that even thin leaves around a fine feature train on
enough points to solve their local problem. The loss weights $\lambda_r$, $\lambda_{ic}$, and $\lambda_{bc}$ are kept from the root
training, and the interface terms inherit them as well: the spatial-face terms
of \eqref{eq:interface_loss} carry $\lambda_{\Gamma} = \lambda_{bc}$, while the
time-face (local initial condition) terms carry $\lambda_{ic}$.
Examples of interface training samples for inner domain experts are shown in
Figure~\ref{fig:local_expert_training}.

\vspace{0.5em}

\begin{figure}[H]
\centering

\begin{subfigure}{\textwidth}
\centering
\includegraphics[width=\textwidth]
{images/section3_method/local_expert_training_schrodinger.png}
\caption{Schrödinger}
\end{subfigure}
\hfill
\begin{subfigure}{\textwidth}
\centering
\includegraphics[width=\textwidth]
{images/section3_method/local_expert_training_allen_cahn.png}
\caption{Allen--Cahn}
\end{subfigure}

\vspace{0.5em}

\begin{subfigure}{\textwidth}
\centering
\includegraphics[width=\textwidth]
{images/section3_method/local_expert_training_burgers.png}
\caption{Burgers}
\end{subfigure}

\vspace{0.5em}

\caption{
Examples of the localized training problem for individual experts, showing interface samples carrying data from the root prediction 
$u_0$. Panels show an interior expert, an expert adjacent to the initial time, and an expert adjacent to the physical boundary.}
\label{fig:local_expert_training}
\end{figure}

\subsection{Smooth Blending and Global Fine-Tuning}
\label{sec:smoothstep_windows}

The trained leaf experts are combined into a single global approximation
through smooth, compactly supported window functions similarly to FBPINNs~\cite{moseley2021fbpinn}. Two
differences are essential here: the decomposition is not a fixed uniform grid but
based on the adaptive tree of \S\ref{sec:tree_construction}, and the experts are first 
trained independently on local sub-problems before they are joined through a partition of unity and trained further under the global loss.

Let $\{\Omega_j\}_{j=1}^{k}$ be the set of leaf sub-domains obtained from the hierarchical
decision tree decomposition, where 
\[
\Omega\times[0,T]=\cup_{j=1}^k \Omega_j
\]
and each region is an axis-aligned hyperrectangle
\[
\Omega_j = \Bigl(\textstyle\prod_{i=1}^{d} [a_{j,i}, b_{j,i}]\Bigr) \times [a_{j,t}, b_{j,t}].
\]
To smoothly blend the contributions of different experts while
ensuring compatibility with the PDE's differential order, we define compactly
supported window functions with configurable smoothness. Let $N_w$ denote the
window smoothness order, chosen such that $N_w \geq N$, the differential order
of the PDE (e.g., $N_w \geq 2$ for Burgers or Allen--Cahn, $N_w \geq 3$ for
KdV, $N_w \geq 4$ for KS). Since the domain decomposition operates over all dimensions including time, no dimension
receives special treatment and the window must be sufficiently smooth for all
derivatives.

We first define the smoothstep polynomial $S_{N_w} : [0,1] \to [0,1]$,
which satisfies $S_{N_w}(0) = 0$, $S_{N_w}(1) = 1$, and
$S_{N_w}^{(k)}(0) = S_{N_w}^{(k)}(1) = 0$ for $k = 1, \ldots, N_w$:
\begin{align}
S_1(t) &= 3t^2 - 2t^3 \quad \text{($C^1$)} \\
S_2(t) &= 10t^3 - 15t^4 + 6t^5 \quad \text{($C^2$)} \\
S_3(t) &= 35t^4 - 84t^5 + 70t^6 - 20t^7 \quad \text{($C^3$)} \\
S_4(t) &= 126t^5 - 420t^6 + 540t^7 - 315t^8 + 70t^9 \quad \text{($C^4$)}
\end{align}
The one-sided compact ramp function is defined as:
\begin{equation}
\rho_{N_w}(s) = \begin{cases}
0 & \text{if } s \leq 0 \\
S_{N_w}(s) & \text{if } 0 < s < 1 \\
1 & \text{if } s \geq 1
\end{cases}
\end{equation}
For region $\Omega_j$ with bounds $[a_{j,i}, b_{j,i}]$ along
the $i$-th spatial coordinate, $1\le i\le d$, or the time
coordinate $t$, let $\delta_{j,i} = \alpha (b_{j,i} - a_{j,i})$ denote the
overlap (transition) half-width, where $\alpha > 0$ is a configurable fraction. The one dimensional window along the $i$-th coordinate is:
\begin{equation}
\omega_{j,i}(z) = \rho_{N_w}\!\left(\frac{z - (a_{j,i} - \delta_{j,i})}{\delta_{j,i}}\right) \cdot \rho_{N_w}\!\left(\frac{(b_{j,i} + \delta_{j,i}) - z}{\delta_{j,i}}\right)
\label{eq:window_1d}
\end{equation}
This produces a flat-top window: $\omega_{j,i}(z) = 1$ for
$z \in [a_{j,i}, b_{j,i}]$, smooth $C^{N_w}$ ramps in the overlaps
$[a_{j,i} - \delta_{j,i}, a_{j,i}]$ and $[b_{j,i}, b_{j,i} + \delta_{j,i}]$,
and exactly zero outside
$[a_{j,i} - \delta_{j,i}, b_{j,i} + \delta_{j,i}]$.

The region indicator function is therefore the tensor product of the one dimensional windows over all coordinates:
\begin{equation}
\Psi_{\Omega_j}(x,t) = \omega_{j,t}(t) \prod_{i=1}^{d} \omega_{j,i}(x_i),
\qquad x = (x_1,\ldots,x_d).
\label{eq:smoothstep_windows}
\end{equation}

\begin{figure}[H]
\centering

\begin{subfigure}{0.48\textwidth}
    \centering
    \includegraphics[width=\textwidth]
    {images/section3_method/window_1d_slice.png}
\end{subfigure}
\hfill
\begin{subfigure}{0.48\textwidth}
    \centering
    \includegraphics[width=\textwidth]
    {images/section3_method/window_2d_heatmaps.png}
\end{subfigure}

\caption{
Compact smoothstep window functions ($C^1$, $C^2$, $C^4$ and overlap fraction $\alpha = 0.2$) for a
representative domain decomposition. Left: one dimensional slices for several regions, showing the flat-top
response and smooth overlap transitions. Right: two dimensional indicator
heatmaps $\Psi_{\Omega_j}(x,t)$, illustrating the compact support and smooth blending overlaps.
}
\label{fig:window_demo}
\end{figure}

Given the indicators, the global solution of AToE over the set of leaves $\Lambda:=\{\Omega_j \}$ of the trimmed tree is
\begin{equation}
u_\theta(x,t)=
\sum_{\Omega_j \in \Lambda}
\tilde{\Psi}_{\Omega_j}(x,t)\, u_j(x,t),
\label{eq:global_atoe}
\end{equation}
with the  partition of unity 
\begin{equation}
\tilde{\Psi}_{\Omega_j}(x,t):= \frac{\Psi_{\Omega_j}(x,t)}{\sum_{\Omega_k \in \Lambda} \Psi_{\Omega_k}(x,t)}.
\label{eq:normalization_atoe}
\end{equation}
Example normalized blending weights are shown in
Figure~\ref{fig:normalization_weights}.

\begin{figure}[H]
\centering
\includegraphics[width=0.6\textwidth]
{images/section3_method/soft_weights_atoe-leaves.png}
\caption{
Example normalized blending weights for the AToE leaf composition. Dashed
rectangles indicate expert region bounds. The leaf experts form a POU over the
domain, each covering its region with smooth transitions in the overlap zones.
}
\label{fig:normalization_weights}
\end{figure}

Finally, the full composition \eqref{eq:global_atoe} is jointly fine-tuned under the global PDE loss. This last stage reconciles the independently trained experts in the overlap zones,
producing a single, globally consistent solution. The fine-tuning stage consists of 40k SSBroyden iterations with no Adam warm-up for the second order benchmarks, and 30k L-BFGS iterations, likewise with no warm-up, for the higher order ones, the latter with a learning rate of $10^{-3}$, smaller than in the earlier stages, as the composed loss landscape tolerates only smaller line-search trial steps. Fine-tuning uses a sampling scheme that draws a fixed fraction of the collocation points from within the overlap regions; we use an overlap sampling fraction of $0.3$, with an
overlap fraction $\alpha = 0.03$ for all benchmarks.

Concentrating samples in the overlaps is
motivated by two observations. First, these are the regions where the
expert sub-networks trained on the boundary conditions provided by the initial coarse approximation $u_0$. Second, typically, the tree approximation produces overlaps that are located exactly on sharp transitions or shocks of the solution. Thus, increased concentrated sampling at the overlaps produces a loss function that guides the  training process to optimize the error on the difficult regions. This is a solution-driven counterpart to residual-based adaptive sampling methods such as RAD~\cite{wu2023rad}, with the sampling density fixed in advance by the decomposition rather than updated from residual during training.